\section{Conclusion}
\label{sec:conclusion}

This dissertation set out to answer how modern web application
vulnerabilities can be authentically modelled and pedagogically
structured within Capture-the-Flag environments to enhance
learning outcomes in cybersecurity education. The answer advanced
here is a seven-challenge suite that combines broad OWASP
coverage---including under-represented classes such as server-side
template injection, server-side request forgery, and insecure
deserialisation---with a deterministic-HMAC personalisation layer
that makes flag sharing attributable and makes LLM-generated
flags structurally invalid, packaged as self-contained Docker
Compose stacks whose intended exploit paths are verified by an
objective end-to-end test harness.

\subsection{Main findings}

Three findings stand out from the evaluation.

First, \textit{ground-truth testing of exploitation is feasible and
cheap}. Writing a Python script that drives each challenge's
intended exploit chain end-to-end and asserts the presence of the
per-user flag in the response body produces an objective,
repeatable measure of solvability that a conventional unit-test
suite cannot. The scripts serve simultaneously as regression
tests, as evidence for markers, and as the controlled baseline
against which the LLM track is interpreted. This pattern
generalises well beyond the present suite.

Second, \textit{flag-layer personalisation is a surprisingly strong
integrity posture at a fraction of the cost of full environment
randomisation}. The Chothia--Novakovic SecGen
framework~\cite{chothia2023} correctly identifies per-student
randomisation as the strong answer to academic integrity, but the
full cost of SecGen-style VM parameterisation is disproportionate
for a web-specific graded module. Restricting personalisation to
the flag layer (with username embedded in every flag and a
per-deployment salt that invalidates write-ups across years) gives
up the methodology-sharing defence in exchange for a marked
reduction in operational overhead. For web-exploitation coursework,
where methodology sharing generally still requires competence to
act on, this is a pragmatic trade.

Third, \textit{scaffolding design is the main determinant of
pedagogical accessibility in the Advanced tier}. The LLM and human
tracks agreed that the distinguishing feature of a successful
Advanced-tier player was their willingness to engage with the
breadcrumbing artefacts (CHANGELOG hints, \texttt{/health}
metadata, WAF source disclosure) rather than to type a canonical
public payload. This reframes the design question: the question is
not ``how hard is the vulnerability'' but ``how visible is the
reasoning path that solves it.''

\subsection{Contributions}

The specific contributions of this dissertation are:
\begin{itemize}
    \item a seven-challenge web-exploitation suite covering six
          of the ten 2021 OWASP categories, with three distinct
          injection sub-types and three stacks rarely
          represented in educational CTFs (Flask/Jinja2,
          Rust/Actix-web, React with queued admin bot);
    \item a deterministic HMAC-SHA256 flag personalisation
          scheme with per-deployment and per-flag-slot salt
          parameterisation, embedded-username attribution, and
          formal integrity properties at the flag layer;
    \item a container-native distribution model that keeps
          student data strictly out of the image layer, making
          the same image safely reusable across cohorts;
    \item an \textit{exploit-as-a-test} end-to-end harness that
          acts as the ground-truth oracle for challenge
          solvability;
    \item a dual evaluation methodology combining automated
          exploit replay with a structured LLM solvability trial
          and a qualitative human walkthrough; and
    \item documented design rationale, unintended-vulnerability
          audits, and development history for each challenge,
          captured as first-class artefacts in the repository.
\end{itemize}

\subsection{Further work}

Three strands of further work are worth pursuing.

\textit{Statistical evaluation of learning gain.} The human track
here is diagnostic rather than confirmatory. A properly powered
pre-/post-test study, using the normalised-gain
statistic~\cite{hake1998} over a full cohort of Networks-and-Systems
students, would provide the measurement that both
Meinsma et al.~\cite{Meinsma2022} and Vykopal et al.~\cite{Vykopal}
argue the field is missing. Because the suite ships with its own
objective solvability baseline, such a study can separate
``challenge was solvable'' from ``student solved it'' more
cleanly than prior platforms have permitted.

\textit{Expanded LLM agentic harnesses.} The LLM evaluation here
uses two scaffolding conditions (passive and agentic). The space
of agentic harnesses continues to expand~\cite{ji2025} and an
adversarial-evolution style trial---where each iteration of the
suite attempts to resist the best-known harness---would give the
integrity argument a moving-target form that matches the
real-world threat.

\textit{Broader OWASP coverage and environment-level
randomisation.} Four OWASP categories (A04, A06, A08, A09) remain
outside the current suite's coverage, either because they are
difficult to flag-ify (A04 Insecure Design) or because they
require infrastructure outside the web layer (A08, A09). Extending
the personalisation layer toward SecGen-style
environment randomisation for a subset of challenges would close
the methodology-sharing gap without replacing the whole scheme.

\subsection{Closing remark}

The contribution of this project is ultimately a concrete
demonstration that the three commonly-separated goals of CTF
design---vulnerability realism, pedagogical accessibility, and
academic integrity---can be pursued jointly in a single web-specific
coursework context, without either compromising any of them or
paying the full cost of per-student environment randomisation. The
suite produced here is, to the author's knowledge, the first
web-exploitation CTF corpus to combine
deterministic-HMAC personalisation, container-native distribution,
and ground-truth exploit-script testing in an evaluated form, and
it is offered as a reusable baseline for future work in
cybersecurity education.
